\documentclass[11pt]{article}
\usepackage{amsmath,amssymb,amsthm,latexsym}
\usepackage{mathrsfs,mathtools}
\usepackage{graphicx}
\usepackage{comment}
\usepackage{bm}
\usepackage{natbib}
\usepackage{subfigure,multirow}
\newtheorem{theorem}{Theorem}
\newtheorem{lemma}{Lemma}
\usepackage[colorlinks,linkcolor=black,citecolor=black,urlcolor=black]{hyperref}
\usepackage{verbatim}

\renewcommand\baselinestretch{1.25}
\topmargin -.5in
\textheight 9in
\textwidth 6.6in
\oddsidemargin -.25in
\evensidemargin -.25in

\title{Solutions to Homework 6}
\author{Xiaochuan Gong}
\date\today

\begin{document}
\maketitle

\subsection*{Problem 1}
Let $X=$
\begin{pmatrix}
    Y \\
    Z
\end{pmatrix}
$\sim \mathcal{N}(\mu,\Sigma)$ with $\mu=$
\begin{bmatrix}
    \mu_y \\
    \mu_z
\end{bmatrix}
, $\Sigma=$
\begin{pmatrix}
    \Sigma_{yy} & \Sigma_{yz} \\
    \Sigma_{zy} & \Sigma_{zz}
\end{pmatrix}
$\succcurlyeq 0$, then $Z|Y \sim \mathcal{N}(\mu_{z|y},\Sigma_{z|y})$ with
\begin{equation}
    \label{1}
    \begin{aligned}
        & \mu_{z|y} = \mu_z + \Sigma_{zy}\Sigma_{yy}^{-1}(y-\mu_y), \\
        & \Sigma_{z|y} = \Sigma_{zz} - \Sigma_{zy}\Sigma_{yy}^{-1}\Sigma_{yz}.
    \end{aligned}
\end{equation}


\begin{proof}
We first assume that $Y$ is an $n_1\times 1$ vector, $Z$ is an $n_2\times 1$ vector and $X$ is an $n\times 1$ vector, where $n=n_1+n_2$. By construction, the joint distribution of $Y$ and $Z$ is $\mathcal{N}(\mu,\Sigma)$. Moreover, the marginal distribution of $Y$ is $\mathcal{N}(\mu_y,\Sigma_{yy})$. According to the law of conditional probability, it holds that
\[
p(z|y) = \frac{p(z,y)}{p(y)}.
\]
Since we have 
\[
p(z|y) = \frac{\mathcal{N}(x;\mu,\Sigma)}{\mathcal{N}(y;\mu_y,\Sigma_{yy})},
\]
and then we can use the probability density function of the multivariate normal distribution, this becomes:
\begin{equation}
    \label{2}
    \begin{aligned}
        p(z|y) & = \frac{\frac{1}{\sqrt{(2\pi)^{n}|\Sigma|}}\cdot\text{exp}\left[-\frac{1}{2}(x-\mu)^{T}\Sigma^{-1}(x-\mu)\right]}{\frac{1}{\sqrt{(2\pi)^{n_1}|\Sigma_{yy}|}}\cdot\text{exp}\left[-\frac{1}{2}(y-\mu_y)^{T}\Sigma_{yy}^{-1}(y-\mu_y)\right]} \\
        & = \frac{1}{\sqrt{(2\pi)^{n-n_1}}}\cdot\sqrt{\frac{|\Sigma_{yy}|}{|\Sigma|}}\cdot\text{exp}\left[-\frac{1}{2}(x-\mu)^{T}\Sigma^{-1}(x-\mu)+\frac{1}{2}(y-\mu_y)^{T}\Sigma_{yy}^{-1}(y-\mu_y)\right].
    \end{aligned}
\end{equation}
Write the inverse of $\Sigma$ as 
\[
\Sigma^{-1} = 
\begin{bmatrix}
    (\Sigma_{yy}-\Sigma_{yz}\Sigma_{zz}^{-1}\Sigma_{zy})^{-1} & -(\Sigma_{yy}-\Sigma_{yz}\Sigma_{zz}^{-1}\Sigma_{zy})^{-1}\Sigma_{yz}\Sigma_{zz}^{-1} \\
    -\Sigma_{zz}^{-1}\Sigma_{zy}(\Sigma_{yy}-\Sigma_{yz}\Sigma_{zz}^{-1}\Sigma_{zy})^{-1} & \Sigma_{zz}^{-1}+\Sigma_{zz}^{-1}\Sigma_{zy}(\Sigma_{yy}-\Sigma_{yz}\Sigma_{zz}^{-1}\Sigma_{zy})^{-1}\Sigma_{yz}\Sigma_{zz}^{-1} \\
\end{bmatrix}
\triangleq
\begin{bmatrix}
    \Sigma^{11} & \Sigma^{12} \\
    \Sigma^{21} & \Sigma^{22}
\end{bmatrix}
\]
and plugging this into \eqref{2}, we have 
\begin{equation*}
    \begin{aligned}
        p(z|y) & = \frac{1}{\sqrt{(2\pi)^{n-n_1}}}\cdot\sqrt{\frac{|\Sigma_{yy}|}{|\Sigma|}}\cdot\text{exp}\left[-\frac{1}{2}(x-\mu)^{T}\Sigma^{-1}(x-\mu)+\frac{1}{2}(y-\mu_y)^{T}\Sigma_{yy}^{-1}(y-\mu_y)\right] \\
        & = 
        \begin{split}
            & \frac{1}{\sqrt{(2\pi)^{n-n_1}}}\cdot\sqrt{\frac{|\Sigma_{yy}|}{|\Sigma|}}\cdot\text{exp}\bigg\{-\frac{1}{2}\left[(y-\mu_y)^{T}\Sigma^{11}(y-\mu_y)+2(y-\mu_y)^{T}\Sigma^{12}(z-\mu_z) \\
            & +(z-\mu_z)^{T}\Sigma^{22}(z-\mu_z)\right]+\frac{1}{2}(y-\mu_y)^{T}\Sigma_{yy}^{-1}(y-\mu_y)\bigg\}
        \end{split} \\
        & = 
        \begin{split}
            & \frac{1}{\sqrt{(2\pi)^{n-n_1}}}\cdot\sqrt{\frac{|\Sigma_{yy}|}{|\Sigma|}}\cdot\text{exp}\bigg\{-\frac{1}{2}\left[z-(\mu_z+\Sigma_{zy}^{T}\Sigma_{yy}^{-1}(y-\mu_y))\right]^{T}(\Sigma_{zz}-\Sigma_{zy}\Sigma_{yy}^{-1}\Sigma_{yz})^{-1} \\
            & \left[z-(\mu_z+\Sigma_{zy}^{T}\Sigma_{yy}^{-1}(y-\mu_y))\right]\bigg\}, 
        \end{split}
    \end{aligned}
\end{equation*}
where we use the fact that $\Sigma_{zy}^{T}=\Sigma_{yz}$ and that ${\Sigma^{21}}^{T}=\Sigma^{12}$ since both $\Sigma$ and $\Sigma^{-1}$ are symmetric matrices.
Note that the determinant of $\Sigma$ is
\[
|\Sigma| = 
\begin{vmatrix}
    \Sigma_{yy} & \Sigma_{yz} \\
    \Sigma_{zy} & \Sigma_{zz} 
\end{vmatrix}
= |\Sigma_{yy}|\cdot|\Sigma_{zz}-\Sigma_{zy}\Sigma_{yy}^{-1}\Sigma_{yz}|.
\]
With this and $n=n_1+n_2$, we finally arrive at
\[
\begin{split}
    p(z|y) = 
    & \frac{1}{\sqrt{(2\pi)^{n_2}|\Sigma_{zz}-\Sigma_{zy}\Sigma_{yy}^{-1}\Sigma_{yz}|}}\cdot\text{exp}\bigg\{-\frac{1}{2}\left[z-(\mu_z+\Sigma_{zy}^{T}\Sigma_{yy}^{-1}(y-\mu_y))\right]^{T}(\Sigma_{zz}-\Sigma_{zy}\Sigma_{yy}^{-1}\Sigma_{yz})^{-1} \\
    & \left[z-(\mu_z+\Sigma_{zy}^{T}\Sigma_{yy}^{-1}(y-\mu_y))\right]\bigg\}, 
\end{split}
\]
which is exactly the density function of a multivariate normal distribution $\mathcal{N}(\mu_{z|y},\Sigma_{z|y})$. 
\end{proof}
\end{document}